\section{The literal constant-phase equivalence}
\label{sec:dichotomy}

One soluble exact-content key has
\[
 b=c\kappa,\qquad
 B=\frac{b}{(b,\sigma)},
\]
and phase slope
\begin{equation}
 \Omega=\ell v\sigma B.
\label{eq:phase-slope}
\end{equation}
Literal target provenance gives \(c\mid\Omega\), and TPC-94 therefore
reduces the algebraic phase conductor to \(1\) or \(q_X\).  For every
nonzero low frequency, conductor \(1\) is exactly \(q_X\mid\Omega\).

\begin{theorem}[Literal constant-phase equivalence]
\label{thm:constant-dichotomy}
On every opened-active literal key,
\begin{equation}
 q_X\mid\Omega
 \quad\Longleftrightarrow\quad
 q_X\mid\sigma.
\label{eq:constant-equivalence}
\end{equation}
\end{theorem}

\begin{proof}
\Cref{lem:content-step-below-q} gives \(0<B<q_X\).
The prime \(q_X\) divides neither \(\ell\) nor \(v\) by
\cref{lem:padding-consequences}.  Euclid's lemma applied to
\(\Omega=\ell v\sigma B\) gives the forward implication in
\eqref{eq:constant-equivalence}; the reverse implication is
immediate.
\end{proof}

\begin{corollary}[No hidden progression-step alternative]
\label{cor:qB-return}
The general algebraic possibility \(q_X\mid B\) is empty on
opened-active literal support.  Thus every opened-active
conductor-one key, including a one-point key, retains the target
divisibility \(q_X\mid\sigma\).
\end{corollary}

\begin{remark}[Algebraic conductor versus interval constancy]
On a singleton interval every character is pointwise constant, even
when its algebraic conductor is \(q_X\).  The sector treated here is
the conductor-one sector \(q_X\mid\Omega\), not the larger class of
all characters whose restriction happens to be constant on one
point.  Those nonconstant-conductor singletons remain in the short
sector.
\end{remark}

\begin{remark}[Why primality matters]
For a composite padding modulus, its prime factors could split among
\(\ell,v,\sigma,B\), and \eqref{eq:constant-equivalence} would
not follow.  The prime no-wrap choice is doing genuine arithmetic
work here, although it remains an auxiliary embedding.
\end{remark}
